% Глава 2. Материал: docs/pedagogy/, docs/research/literature-review.md, docs/research/linters-hooks-agents.md,
% reports/ (SRC-038). Таблица Т-09 (docs/thesis/figures.md).
\chapter{Обзор литературы}\label{ch:review}

Глава рассматривает шесть областей, с которыми должен считаться тренажёр: букварь и его метод
(раздел~\ref{sec:lr-primer}), учеников и их трудности (раздел~\ref{sec:lr-learners}), исследования
компьютерного обучения произношению (раздел~\ref{sec:lr-capt}), автоматическую оценку чтения Корана
(раздел~\ref{sec:lr-mdd}), методы адаптации речевой модели к новой аудитории (раздел~\ref{sec:lr-adapt})
и программную инженерию с ИИ-агентами (раздел~\ref{sec:lr-agents}).
Раздел~\ref{sec:lr-related} сравнивает существующие инструменты, раздел~\ref{sec:lr-gap} подводит итог.
Обзор нарративный и обзорный (scoping); поиск выполнен в октябре 2026~г., у каждого источника есть
оценка надёжности в реестре проекта \cite{SRC-038}.

\section{Букварь и его метод}\label{sec:lr-primer}

\subsection{Максуди и два букваря}

Ахмад-Хади Максуди (1868--1941), казанский преподаватель арабского языка и педагог-джадид, написал два
букваря, которые часто путают. «Мөгаллим әүвәл» («Первый учитель», 1892) учил татарской грамоте на
арабской графике и выдержал более тридцати изданий до 1918~г. \cite{SRC-015,SRC-019}. «Муаллим ас-Сани»
(«Второй учитель») продолжал его и вёл от букв к правильному чтению Корана с основами \emph{таджвида}.
Его обычно относят к середине 1890-х, но первичного библиографического подтверждения первого издания не
найдено \cite{SRC-018}. Современные тексты восходят к казанскому изданию 1917~г., на котором основано и
пособие для русскоязычных \cite{SRC-024}. В 2014~г. Духовное управление мусульман Казахстана запустило
проект распространения книги по мечетям и медресе \cite{SRC-020}. В работе используется альметьевское
издание 2010~г., выпущенное при Духовном управлении мусульман Татарстана в орфографии Мединского мусхафа
\cite{SRC-040}.

\subsection{Звуковой метод}

Татарская школа того времени различала три метода обучения грамоте \cite{SRC-017}: буквослагательный
(\emph{усул-и хәрфия}), при котором сначала учили названия букв; слоговой (\emph{усул-и мәддия}); и
звуковой (\emph{усул-и саутия}) реформаторов, при котором буква изучается как звук, а слоги собираются
осознанно. Переходу к звуковому методу приписывают сокращение обучения чтению «с нескольких лет до одного
учебного года» \cite{SRC-016}. Дамиров отмечает, что реформатор З.~Расулев обосновывал метод традицией
\emph{таджвида}, а не европейской педагогикой \cite{SRC-016}, --- довод, важный для религиозной аудитории:
звуковой метод предстаёт как внутренний для исламской традиции.

Фахрутдинова и Ибрагимов описывают вклад Максуди как отказ от заучивания названий букв, фонетическую
ориентацию на артикуляцию и звучание, постепенное и осмысленное продвижение; его пятнадцать дидактических
принципов 1906~г. сводятся к четырём группам: постепенность и преемственность; индивидуализация и
активность; наглядность и апперцепция; практика и непрерывный контроль понимания \cite{SRC-015}.
Таблица~\ref{tab:lr-principles} показывает, как каждая группа превращается в требование к цифровому
тренажёру; та же таблица направляла решения по интерфейсу в главе~\ref{ch:impl}.

\begin{table}[ht]
\centering\small
\caption{Принципы Максуди \cite{SRC-015} и их воплощение в функциях тренажёра.}
\label{tab:lr-principles}
\begin{tabularx}{\linewidth}{@{}P{4.4cm}L@{}}
\toprule
Принцип & Функция тренажёра \\
\midrule
Постепенность и преемственность & Уроки в порядке книги; строка освоена, когда прочитана верно; урок не
длиннее дюжины строк. \\
Индивидуализация и активность & Свой темп и история; повтор слабых строк; сколько угодно попыток. \\
Наглядность и апперцепция & Эталонное аудио до чтения; буква урока красным, как в книге; подсказки,
связывающие новые звуки со знакомыми. \\
Практика и непрерывный контроль & Каждая попытка проверяется сразу; каждая ошибка объясняется; сразу
можно повторить. \\
Гуманность & Без штрафов и рейтингов; серия дней как поощрение. \\
Звук прежде названия & Карточка учит звук с огласовкой; название буквы --- только справка. \\
\bottomrule
\end{tabularx}
\end{table}

\subsection{Строение книги и цикл урока}

Книга не следует алфавиту. Два независимых описания и издание 2010~г. сходятся: буквы вводятся
группами, начиная с хамзы и \ar{ر ز م} --- букв простой, в основном несоединяемой формы, из которых сразу
складываются слоги вроде \ar{زَرَ} и \ar{مَرَ} \cite{SRC-021,SRC-040}. За буквами идут долгие гласные,
\emph{шадда} и \emph{танвин}; хамза, артикль и \emph{сукун}; правила \emph{мим} и \emph{нун}; особенности
письма Корана; остановки и долготы чтения; и, наконец, короткие суры. Каждая группа букв заканчивается
строками «обратите внимание», где противопоставлены похожие звуки.

Видеокурсы и занятия в мечетях по этой книге следуют одному циклу урока: послушать образец, повторить
его (в одном курсе --- каждую строку трижды), прочитать самостоятельно, получить исправление и идти
дальше только после безошибочного чтения \cite{SRC-021,SRC-022}. Это цифровая форма \emph{таляккы} ---
устной передачи, при которой учитель читает, ученик повторяет, а учитель сразу поправляет
\cite{SRC-058}. Индонезийский метод Икро делает роль учителя явной: когда буквы выучены, учитель больше
не объясняет, а показывает, где ошибка, и ученик исправляет её сам \cite{SRC-057}. Преподаватели также
советуют не опираться на транскрипцию, потому что «на бумаге невозможно передать звук»
\cite{SRC-021,SRC-022}, а после курса нормой остаётся еженедельная устная проверка у учителя
\cite{SRC-039}. Контролируемых испытаний букварных методик чтения Корана --- «Муаллим ас-Сани», «Нурани
Каида» или Икро --- не найдено \cite{SRC-038}.

\section{Ученики}\label{sec:lr-learners}

Целевые ученики --- русскоязычные взрослые в России и Казахстане, которые учатся читать по-арабски с нуля.
Рецензируемых исследований их ошибок чтения нет. Методические источники для русскоязычных перечисляют
ожидаемые трудности: фарингальные \ar{ح} и \ar{ع}; увулярный \ar{ق}, произносимый как [к]; смешение \ar{غ}
и \ar{خ}; межзубные \ar{ث}, \ar{ذ}, \ar{ظ}; эмфатические \ar{ص}, \ar{ض}, \ar{ط}, \ar{ظ} без твёрдости;
укороченные долгие гласные; неудвоенная \emph{шадда}; \ar{ر}, смягчённый перед гласными переднего ряда
\cite{SRC-023}. Исследования носителей других языков сходятся в том, что фарингальные, альвеолярно-дентальные
и межзубные звуки труднее всего для неарабов \cite{SRC-034}, а индонезийцы больше всего ошибаются на
\ar{ع} и \ar{ض} \cite{SRC-035}. Собственные минимальные пары книги --- \ar{ق}/\ar{ك}, \ar{ث}/\ar{س},
\ar{ص}/\ar{س}, \ar{ط}/\ar{ت}, \ar{ح}/\ar{خ}, \ar{ع}/\ar{غ}, \ar{ض}/\ar{د}, \ar{ذ}/\ar{ز} --- нацелены
на те же противопоставления \cite{SRC-040}.

Фонология подсказывает одно различие между группами: в казахском есть увулярные [q] и [\ipa{ʁ}], поэтому
\ar{ق} и \ar{غ} казахоязычным, возможно, даются легче, чем русскоязычным. Это гипотеза (разведочная E-1), а не
вывод. Практическая транскрипция для русскоязычных тоже требует внимания: обычная арабско-русская
практическая транскрипция теряет различие \ar{ص} и \ar{س}, а научная система отмечает эмфатические и
фарингальные точкой снизу, а межзубные --- чертой снизу \cite{SRC-051}.

\section{Компьютерное обучение произношению}\label{sec:lr-capt}

\subsection{Эффективность}

Метаанализ Нго, Чена и Лая об автоматическом распознавании речи в обучении английскому произношению
показал средний общий эффект ($g=0{,}69$). Эффект больше для сегментных признаков, чем для
сверхсегментных, и больше при явной корректирующей обратной связи, чем при неявной, например простом
транскрипте \cite{SRC-025}. Обучение произношению в целом даёт большой эффект \cite{SRC-105}, так что
автоматическую обратную связь правильнее считать способом дать больше обучения, а не лучшее обучение.
Систематический обзор Амрате и Цай добавляет, что распознаватели ошибаются на речи с акцентом и детской
речи и дают ложные замечания на понятное произношение, и рекомендует использование под контролем, когда
учитель или партнёр может поправить систему \cite{SRC-026}.

Важна и тренировка восприятия. Тренировка на голосах многих дикторов (HVPT) даёт средний или большой и
устойчивый эффект на различение звуков \cite{SRC-044} и меньший, но реальный перенос на произношение
\cite{SRC-045}. Для запоминания распределённая практика --- один из самых надёжных эффектов науки об
обучении \cite{SRC-107}. Тренажёр, который проверяет только произношение, оставляет место для
упражнений на слух и интервального повторения; проект записывает их как будущие функции.

\subsection{Какой должна быть обратная связь}

Листер и Ранта обнаружили, что рекаст --- учитель просто произносит правильную форму --- редко приводит к
самоисправлению, а подсказки, указывающие место ошибки и предлагающие попробовать снова, приводят чаще
\cite{SRC-053}. Нери, Куккьярини и Стрик считают, что обратная связь по произношению должна касаться
немногих важных ошибок, а не всех, а их поздняя работа показала, что автоматическая обратная связь
исправляет именно разобранные ошибки \cite{SRC-054}. Обзор Шют о формирующей обратной связи советует
конкретность, простоту и небольшие порции с подробностями по запросу \cite{SRC-055}. Хэтти и Тимперли
называют самым сильным компонентом ответ на вопрос «что делать дальше» и находят похвалу личности почти
бесполезной \cite{SRC-056}. Коранская педагогика упорядочивает ошибки сходно: явные ошибки, меняющие букву
или огласовку (\emph{лахн джали}), исправляют раньше скрытых, вроде долготы или назализации (\emph{лахн
хафи}) \cite{SRC-052}.

Эти выводы сходятся к единице обратной связи из трёх частей: \emph{что} не так и где, \emph{как} сказать и
\emph{что делать дальше} --- для небольшого числа ошибок за раз, явные ошибки первыми. Коммерческие
приложения добавляют четвёртый элемент --- подсветку ошибки прямо в тексте \cite{SRC-046}.

\subsection{Мессенджеры как среда}

Telegram --- повседневный мессенджер целевой аудитории, голосовые сообщения в нём встроены. Четырёхнедельное
исследование показало, что Telegram-канал улучшил английское произношение иранских студентов
\cite{SRC-092}, но контролируемых исследований Telegram-бота с автоматической проверкой произношения не
найдено. Правила платформы разрешают собирать данные для машинного обучения только с явного, активного и
отзываемого согласия \cite{SRC-033}; лимиты на сообщения и файлы тоже ограничивают дизайн \cite{SRC-060}.

\section{Автоматическая оценка чтения Корана}\label{sec:lr-mdd}

\subsection{Задача и метрики}

Обнаружение и диагностика ошибок произношения (MDD) сравнивают фонемы ученика с каноническими фонемами
текста. Стандартный протокол, использованный в L2-ARCTIC \cite{SRC-063} и принятый в коранических
бенчмарках \cite{SRC-064,SRC-066}, выравнивает три последовательности --- каноническую, размеченную и
распознанную --- и относит каждую каноническую фонему к верному принятию (TA), верному отклонению (TR,
настоящая ошибка найдена), ложному принятию (FA, ошибка пропущена) или ложному отклонению (FR, ложная
тревога). Отсюда precision, recall и F1:
\begin{equation}
  P = \frac{TR}{TR+FR},\qquad R = \frac{TR}{TR+FA},\qquad F_1 = \frac{2PR}{P+R}.
\end{equation}
Точность диагноза дополнительно проверяет, совпадает ли распознанная замена с размеченной. Для тренажёра
цена двух видов ошибок разная: пропуск замедляет обучение, ложная тревога на верном чтении подрывает
доверие.

\subsection{Бенчмарки}

QuranMB.v1, представленный на Interspeech 2025, стал первым открытым бенчмарком обнаружения ошибок
произношения арабского на материале Корана; его базовые системы имели высокую полноту и очень низкую
точность \cite{SRC-004}. Соревнование Iqra'Eval на ArabicNLP 2025 дало общую оценку, системы участников
достигли F1 около 0,40--0,44 \cite{SRC-043}. В IQRA 2026 на Interspeech, на QuranMB.v2, участвовали
девятнадцать команд; лучшая система получила F1 0,72 по фонемам против 0,44 у базовой
\cite{SRC-031,SRC-086}. Сильнейшие системы --- CTC-распознаватели фонем на самообучаемых энкодерах.
Организаторы отмечают, что в данных нет изучающих арабский как второй язык, и называют их сбор
приоритетом \cite{SRC-031}.

\subsection{Quran Muaalem}

Абдельфаттах, Халил и Аббас построили Quran Muaalem вокруг почти автоматического конвейера данных: 848
часов и около 286 тыс. фрагментов от 22 чтецов-экспертов, размеченных в специальной фонетической записи
Quran Phonetic Script, которая кодирует и фонемы, и артикуляционные признаки (\emph{сифат})
\cite{SRC-002}. Энкодер Wav2Vec2-BERT питает одиннадцать CTC-голов: фонемы и десять признаков, таких как
глухость и звонкость, твёрдость и мягкость, \emph{калькаля}. Модель, код и данные открыты \cite{SRC-003};
пакет отдаёт вероятности по фонемам и признакам, которые тренажёр может использовать для решений с учётом
уверенности \cite{SRC-062}.

Заявленные результаты показывают разрыв доменов. На тесте авторов из чтения экспертов доля ошибок по
фонемам --- 0,54\,\% (0,21\,\% в среднем по всем уровням). На qdat\_bench, 159 записях с настоящими
ошибками, она вырастает до 5,78\,\%, средний F1 по таджвиду --- 0,758, а \emph{ихфа} --- лишь 0,453, потому
что акустически близка к ясному \emph{нуну} \cite{SRC-002}. Авторы отмечают, что в обучающих данных только
мужчины-эксперты и только риваят Хафс, а qdat\_bench размечен одним человеком на одном аяте. Для букваря
следуют ещё два ограничения: фонетизатор рассчитан на текст Корана, поэтому отдельным слогам нужны
канонические строки фонем, а медленное послоговое чтение начинающих лежит вне распределения обучающих
данных.

\subsection{Данные учеников и смежные работы}

Ближайший открытый набор чтения неарабов --- около семи тысяч краудсорсинговых записей от носителей более
чем из одиннадцати стран --- собран в Университете Иннополис \cite{SRC-032}; носителей русского и
казахского в нём нет. Работы той же группы --- ближайшие соседи этой диссертации. В одной ИИ-агенты
улучшают проверку заучивания по транскриптам Telegram-бота, данные и оценщик заморожены хешами SHA-256,
сформулированы правила против утечек и подгонки метрики \cite{SRC-067,SRC-069}. В другой предложена
таксономия событий чтения, отделяющая неисправленные ошибки от самоисправлений и безобидных событий вроде
повторов \cite{SRC-068}. Обе проверяют, \emph{что} прочитано --- какие аяты, с пропусками или вставками ---
по транскриптам без огласовок и выносят огласовки и таджвид за рамки. Тренажёр букваря должен проверять,
\emph{как} произнесён каждый слог. Таксономия и правила оценки при этом приняты в работе: самоисправления
размечаются отдельно, а оценщик защищён от ИИ-агента.

\section{Адаптация речевой модели к новой аудитории}\label{sec:lr-adapt}

\subsection{Дообучение на речи изучающих язык}

Модели обнаружения ошибок на самообучаемых речевых энкодерах обычно получают дообучением с CTC на
размеченной речи изучающих язык. Сюй и соавторы показали на L2-ARCTIC, что дообученный wav2vec~2.0 лучше
прежних систем и сохраняет качество, когда размеченных данных мало \cite{SRC-116}; Muaalem построена так
же, только дообучалась на экспертах \cite{SRC-002}. Для нескольких часов речи полное дообучение
600-миллионной модели рискованно: оно переобучается и может испортить то, что модель уже умеет. Метод
LoRA замораживает исходные веса и обучает малые матрицы низкого ранга, добавленные к слоям внимания;
обучаемых параметров на порядки меньше, а исходную модель можно вернуть, просто отключив добавки
\cite{SRC-111}. Дообучение на новой задаче может стирать старые умения --- катастрофическое забывание
\cite{SRC-118}, --- поэтому адаптированную модель нужно проверять и на прежнем домене, здесь --- на чтении
экспертов.

\subsection{Где взять ошибки}

У начинающих ошибок больше, чем у экспертов, но для обучения их всё равно мало, а отдельные пары звуков
встречаются редко. В литературе по MDD данные пополняют синтетическими ошибками: в эталонной
транскрипции заменяют фонемы, как будто ученик ошибся \cite{SRC-119}; стандартный приём распознавания
речи --- изменение скорости записи \cite{SRC-120}. Букварь даёт более естественный источник: его
минимальные пары (\ar{ق}/\ar{ك}, \ar{ص}/\ar{س}, \ar{ح}/\ar{خ} и др.) означают, что верное чтение одной
карточки --- настоящая запись настоящей ошибки для парной. Такое пополнение использует аудио людей, а не
синтез: синтетический арабский голос модель распознаёт плохо (раздел~\ref{sec:eval-prelim}).

\subsection{Уверенность, калибровка и «не уверен»}

Второй путь --- не менять модель, а научиться ей не доверять там, где она ошибается. Классическая мера
Goodness of Pronunciation превращает апостериорные вероятности фонем в оценку произношения
\cite{SRC-115}; поверх таких признаков можно обучить небольшой верификатор на сотнях размеченных примеров.
Нейросети, однако, переуверены, и их вероятности нужно калибровать; стандартная мера --- ожидаемая ошибка
калибровки \cite{SRC-112}. Выборочная классификация разрешает системе отказаться от ответа и описывает
компромисс кривой «риск --- покрытие» \cite{SRC-117}. Конформное предсказание даёт гарантию без допущений
о распределении при перестановочности данных \cite{SRC-113}, а конформный контроль риска позволяет
выбрать порог так, чтобы ожидаемая доля ложных отклонений не превышала заданный уровень \cite{SRC-114}.
Для тренажёра это прямой ответ на главную опасность из раздела~\ref{sec:lr-capt}: ложное «не засчитано» на
верном чтении.

\subsection{Как не обмануть себя}

С малыми данными легко получить завышенный результат. Подбор гиперпараметров на тех же данных, на которых
сообщается качество, смещает оценку вверх; нужна вложенная кросс-валидация или отдельный тест
\cite{SRC-121}. Капур и Нараянан нашли утечку данных в сотнях ML-исследований и предложили её таксономию
\cite{SRC-122}; для речи главная утечка --- один и тот же диктор в обучении и проверке. Отсюда правила
работы: разбиение по ученикам, тест, замороженный до обучения, слепая разметка теста и один прогон на нём.
Те же требования группа научного руководителя сформулировала для исследований, где модели подбирают
ИИ-агенты \cite{SRC-067}. Набор, который будет опубликован, сопровождается карточкой по Гебру и
соавторам \cite{SRC-123}.

\section{Аналоги}\label{sec:lr-related}

Таблица~\ref{tab:T-09} сравнивает существующие инструменты с тренажёром этой работы по признакам,
важным для начинающих букварного уровня.

\begin{table}[ht]
\centering\small
\caption{Аналоги (Т-09). «Проверка» --- автоматическая проверка голоса ученика.}
\label{tab:T-09}
\begin{tabularx}{\linewidth}{@{}P{3.4cm}P{2.0cm}P{2.1cm}L@{}}
\toprule
Система & Уровень & Проверка & Примечания \\
\midrule
Tarteel \cite{SRC-046} & аяты & слова & Заучивание; ошибки подсвечены в тексте; проверки на уровне букв нет. \\
Приложение Muallimi Soniy \cite{SRC-006} & букварь & --- & Порядок правил по книге; узбекская
аудитория. \\
Платформа muallimi\_soniy \cite{SRC-005} & букварь & --- & Открытая веб-версия одного издания. \\
Курс «thani» \cite{SRC-070} & букварь & --- & Тот же порядок букв; тренирует слух (диктанты). \\
Бот Recite Quran \cite{SRC-071,SRC-067} & аяты & транскрипт & Проверка заучивания через распознавание
речи, без огласовок. \\
\textbf{Эта работа} & букварь & фонемы & Порядок книги со страницей и строкой; ошибки на уровне фонем с
объяснением по-русски; полный след решения; разметка учителями. \\
\bottomrule
\end{tabularx}
\end{table}

Инструмента на русском, казахском или узбекском, который проверял бы голосом чтение букварного уровня,
не найдено. Несколько новых приложений обещают ИИ-проверку уровня \emph{Каиды}, но эти обещания взяты со
страниц магазинов и независимо не проверены \cite{SRC-038}.

\section{Программная инженерия с ИИ-агентами}\label{sec:lr-agents}

\subsection{Агенты и обвязка}

ИИ-агент чередует рассуждение, использование инструмента и наблюдение \cite{SRC-108}. SWE-agent показал,
что интерфейс между агентом и компьютером определяет результат: линтер, отклоняющий синтаксически неверные
правки в момент правки, увеличил долю решённых задач в репозиториях \cite{SRC-109}. Практики называют
инфраструктуру вокруг цикла \emph{обвязкой} (harness). Рамка Бёкелер разделяет \emph{направляющие}
(guides), задающие курс до действия (инструкции, скиллы), и \emph{датчики} (sensors), проверяющие результат
(тесты, линтеры, ревьюеры), и различает детерминированные вычислительные проверки и вероятностные
инференциальные; лучше всего обвязка работает для сопровождаемости и архитектурной пригодности, хуже
всего --- для поведения \cite{SRC-029}. Инженерные заметки Anthropic советуют внешние файлы прогресса, одну
функцию за сессию, сквозную проверку и явный запрет править или удалять тесты \cite{SRC-028}. Помощь
LLM в проектировании архитектуры устроена так же: решения разумны, но архитектурные драйверы покрываются
лишь частично, нужен контроль человека \cite{SRC-050}. Разведочное исследование 2\,853 репозиториев
показало, что в практике преобладают файлы контекста, а хуки, скиллы и субагенты редки \cite{SRC-073}.
Рецензируемого определения «harness engineering» пока нет, и источники практиков используются здесь для
понятий, а не для чисел.

\subsection{Данные}

\paragraph{Агенты подгоняют оценку} В ImpossibleBench спецификация противоречит тестам, так что
прохождение означает, что агент воспользовался тестами; передовые модели делали это часто, и у моделей
Claude большая часть таких попыток меняла сами тесты. Тесты только для чтения закрывали этот путь
\cite{SRC-027}. Подгонка под награду, выученная на задачах программирования при обучении, может
обобщаться в более широкое нежелательное поведение, вплоть до саботажа, когда модель потом работает
агентом-программистом \cite{SRC-082}.

\paragraph{Письменные правила слабы} Файлы контекста уровня репозитория вроде \texttt{AGENTS.md} не
повышали успешность задач, иногда снижали её и увеличивали стоимость примерно на пятую часть
\cite{SRC-074}; другое исследование нашло меньшее время и число токенов при той же успешности
\cite{SRC-075}; абляция с двумя агентами и 288 прогонами не нашла влияния на корректность \cite{SRC-077}.
Факторное исследование 1\,650 сессий Claude Code показало, что соблюдение инструкций не зависит от
структуры файла, но падает по ходу сессии \cite{SRC-076}.

\paragraph{Механическая обратная связь помогает} Возврат замечаний статического анализа модели снизил долю
кода с проблемами безопасности с более чем 40\,\% до 13\,\% за десять итераций \cite{SRC-078}. Агенты
способны выводить свойства и писать тесты свойств, находящие реальные ошибки \cite{SRC-080,SRC-081}, а
генерация тестов на основе мутантов в Meta дала тесты, которые инженеры в большинстве случаев приняли
\cite{SRC-079}.

\paragraph{Ощущение --- не измерение} В рандомизированном исследовании опытные разработчики открытого ПО
работали с ИИ-инструментами на 19\,\% дольше, считая, что ускорились \cite{SRC-030}. Агенты, симулирующие
учеников, систематически расходятся с настоящими учениками \cite{SRC-083,SRC-084,SRC-085}, так что они
могут дополнить пилот, но не заменить его.

Вместе эти данные обосновывают принцип, применённый в работе: важные правила должны стать проверками,
которые обеспечивает среда, и каждое срабатывание должно оставлять след, который можно измерить.

\section{Итоги и пробел}\label{sec:lr-gap}

Метод «Муаллим ас-Сани» --- звуковой, пошаговый, построенный на немедленном исправлении --- хорошо
сочетается с автоматической обратной связью, а условия, при которых обратная связь по произношению работает
лучше всего, совпадают с условиями букваря. Открытая модель распознавания фонем существует. Не хватает
доказательств по четырём пунктам:
\begin{enumerate}
  \item насколько точно готовая модель оценивает букварное чтение начинающих, для которых арабский не
    родной, по сравнению с учителями, как часто отклоняет верное чтение и какие из найденных ошибок
    принадлежат модели, а не ученику (RQ-1);
  \item можно ли на малом наборе этой аудитории адаптировать модель так, чтобы ложных отклонений стало
    меньше, а согласие с учителем выше, и не потерять чтение экспертов --- при оценке без утечки (RQ-2);
  \item улучшают ли объяснения на языке ученика обучение на уровне арабских букв (RQ-3);
  \item как ведёт себя обвязка ИИ-агента на протяжении целого исследовательского проекта, включая
    эксперименты обучения (RQ-4).
\end{enumerate}
Следующая глава описывает, как работа собирает эти доказательства.
